%------------------------------------------------------------------------------%

\usepackage{integracao_e_series}

\title{Séries Numéricas -- Teste da Comparação}

%------------------------------------------------------------------------------%
\begin{document}

\addtitle

%------------------------------------------------------------------------------%
%------------------------------------------------------------------------------%
\section{Teste da Comparação}

%------------------------------------------------------------------------------%
\definecolor{c1}{rgb}{0,.7,0}
\definecolor{c2}{rgb}{0,0,.7}
\definecolor{c3}{rgb}{.7,0,0}
\begin{frame}{Teste da Comparação}
  Sejam três séries com termos não negativos
  \[
    {\color{c1} \sum_{n=1}^{\infty} p_n} \qquad
    {\color{c2} \sum_{n=1}^{\infty} a_n} \qquad
    {\color{c3} \sum_{n=1}^{\infty} q_n}
  \]
  com \( \quad {\color{c1}p_n} \;\leq\; {\color{c2}a_n} \;\leq\; {\color{c3}q_n} \quad \) para todo $n>N$

  \pause
  \hspace{\baselineskip}
  \begin{enumerate}
    \item se \tabto{8mm} \(\color{c3} \sum q_n \) convergir \tabto{40mm} então  \(\color{c2} \sum a_n \) converge \\[6mm]\pause
    \item se \tabto{8mm} \(\color{c1} \sum p_n \) divergir  \tabto{40mm} então  \(\color{c2} \sum a_n \) diverge
  \end{enumerate}
\end{frame}

%------------------------------------------------------------------------------%
\addtocounter{exemplo}{1}
\begin{frame}{Exemplo \theexemplo}
  Para analisar a série do \emph{inverso do fatorial}
  \[
    \sum_{n=1}^{\infty} \frac{1}{\,n!\,}
    = 1             \;+\;
      \frac{1}{2!}  \;+\;
      \frac{1}{3!}  \;+\;
      \frac{1}{4!}  \;+\;
      \frac{1}{5!}  \;+\;
      \cdots
  \]
  \pause
  comparamos com a série geométrica
  \[ \sum_{n=1}^{\infty} \frac{1}{2^{n-1}}
    = 1              \;+\;
      \frac{1}{2}    \;+\;
      \frac{1}{2^2}  \;+\;
      \frac{1}{2^3}  \;+\;
      \frac{1}{2^4}  \;+\;
      \cdots
      \qquad
      \qquad
      r=\frac{1}{2}
  \]
  \pause
  A série geométrica converge
  \pause
  e
  \(\quad \frac{1}{n!} < \frac{1}{2^{n-1}}\quad\)
  para $n>2$

  \pause
  \vspace{\baselineskip}
  Então a série do inverso do fatorial converge
\end{frame}

%------------------------------------------------------------------------------%
%------------------------------------------------------------------------------%
\section{Teste da Comparação no Limite}

%------------------------------------------------------------------------------%
\begin{frame}{Teste da Comparação no Limite}
  Sejam as séries
  \(\; \sum {\color{c1} a_n} \;\) e
  \(\; \sum {\color{c3} b_n} \;\)
  com $\;{\color{c1} a_n} > 0\;$ e $\;{\color{c3} b_n}>0\;$ para todo $n>N$
  \pause
  \vspace{\baselineskip}
  \begin{enumerate}
    \item
    \emph{Se} \(\;\lim \frac{\;{\color{c1} a_n}\;}{{\color{c3} b_n}} = c > 0 \) \;\;
    \emph{então}\; ambas convergem ou ambas divergem

    \pause
    \vfill
    \item
    \emph{Se} \(\;\sum {\color{c3} b_n}\) converge
    \tabto{36mm} \emph{e} \(\;\lim \frac{\;{\color{c1} a_n}\;}{{\color{c3} b_n}} = 0\)
    \tabto{66mm} \emph{então} \(\;\sum {\color{c1} a_n} \) também converge

    \pause
    \vfill
    \item
    \emph{Se} \(\;\sum {\color{c3} b_n}\) diverge
    \tabto{36mm} \emph{e} \(\;\lim \frac{\;{\color{c1} a_n}\;}{{\color{c3} b_n}} = \infty\)
    \tabto{66mm} \emph{então} \(\;\sum {\color{c1} a_n} \) também diverge

    \vfill
  \end{enumerate}
\end{frame}

%------------------------------------------------------------------------------%
\addtocounter{exemplo}{1}
\begin{frame}{Exemplo \theexemplo}
  Verificar a convergência da série
  \(
    \sum_{n=1}^\infty \frac{2^n}{3+4^n}
  \)

  \pause
  \vspace{\baselineskip}
  A série converge\pause, pois, tomando
  \[
    b_n
    \;=\; \frac{2^n}{4^n}
    \;=\; \frac{1}{2^n}
  \]
  \pause
  temos
  \[
    \lim \frac{a_n}{b_n}
    \pause
    = \lim \frac{\;\dfrac{2^n}{3+4^n}\;}{\dfrac{1}{2^n}}
    \pause
    = \lim \frac{4^n}{3+4^n}
    \pause
    = \lim \frac{1}{\sfrac{3}{4^n}+1}
    \pause
    = 1
  \]
\end{frame}

%------------------------------------------------------------------------------%
\begin{frame}{Exemplo \theexemplo}
  A série
  \(
    \sum \frac{1}{2^n}
  \)
  converge

  \pause
  \vspace{\baselineskip}
  Pois é uma série geométrica com $\abs{r}=\frac{1}{2}<1$

  \pause
  \vspace{\baselineskip}
  Então a série
  \(
    \sum_{n=1}^\infty \frac{2^n}{3+4^n}
  \)
  também converge
\end{frame}

%------------------------------------------------------------------------------%
%------------------------------------------------------------------------------%
\section{Exemplos}

%------------------------------------------------------------------------------%
\addtocounter{exemplo}{1}
\begin{frame}{Exemplo \theexemplo}
  Use o teste da comparação para mostrar que a série
  \[
    \sum_{n=1}^\infty \frac{n-1}{n^4+2}
  \]
  converge
\end{frame}

%------------------------------------------------------------------------------%
\begin{frame}{Exemplo \theexemplo}
  \begin{align*}
    \onslide<+->{a_n & = \frac{n-1}{n^4+2} \\[3mm]}
    \onslide<+->{& =    \frac{n}{n^4+2} - {\color<3->{gray}\frac{1}{n^4+2}} \\[3mm]}
    \onslide<+->{& \leq \frac{n}{n^4+2}                   \\[3mm]}
    \onslide<+->{& =    \frac{1}{n^3 +  {\color<5->{gray}\sfrac{2}{n}}}      \\[3mm]}
    \onslide<+->{& \leq \frac{1}{n^3}}
  \end{align*}
\end{frame}

%------------------------------------------------------------------------------%
\begin{frame}{Exemplo \theexemplo}
  A série
  \(\;
    \sum_{n=1}^\infty \frac{1}{n^3}
  \;\)
  é uma série $p$ com $p=3>1$, portanto convergente

  \pause
  \vspace{\baselineskip}
  Pelo teste de comparação a série
  \(\;
    \sum_{n=1}^\infty \frac{n-1}{n^4+2}
  \;\)
  também converge
\end{frame}

%------------------------------------------------------------------------------%
\addtocounter{exemplo}{1}
\begin{frame}{Exemplo \theexemplo}
  Use o teste da comparação para verificar a convergência da série
  \(
    \sum_{n=1}^\infty\frac{3^n}{2^n-1}
  \)

  \pause
  \vspace{\baselineskip}
  A série diverge\pause,
  pois
  \[
    a_n = \frac{3^n}{2^n-1}
    \pause
    \;\geq\; \frac{3^3}{2^n}
    \pause
    \;=\;    \left(\frac{3}{2}\right)^{\!\!{n}}
  \]
  \pause
  e a série geométrica
  \(
    \sum_{n=1}^\infty\left(\frac{3}{2}\right)^{\!\!{n}}
  \)
  diverge\pause, pois \quad
  \(
    \abs{r}=\frac{3}{2}>1
  \)
\end{frame}

%------------------------------------------------------------------------------%
\addtocounter{exemplo}{1}
\begin{frame}{Exemplo \theexemplo}
  Utilize o teste da comparação para determinar se a série
  \[
    \sum_{n=1}^{\infty} \frac{\cos^2 n}{n^{\sfrac{3}{2}}}
  \]
  converge ou diverge
\end{frame}

%------------------------------------------------------------------------------%
\begin{frame}{Exemplo \theexemplo}
  Sabemos que \(\abs{\cos x}\leq 1\)\pause,
  portanto
  \[
    0 \;\leq\;\cos^2 x \;\leq\; 1
  \]
  \pause
  como $n^{\sfrac{3}{2}}$ é sempre positivo para $n\geq 1$
  \[
    0 \;\leq\; \frac{\cos^2}{n^{\sfrac{3}{2}}} \;\leq\; \frac{1}{n^{\sfrac{3}{2}}}
  \]
\end{frame}

%------------------------------------------------------------------------------%
\begin{frame}{Exemplo \theexemplo}
  Analisando a série
  \[
    \sum_{n=1}^{\infty} \frac{1}{n^{\sfrac{3}{2}}}
  \]
  percebemos que é uma série $p$ com $p=\frac{3}{2}>1$ e portanto convergente

  \pause
  \vspace{1.5\baselineskip}
  Assim a série original é convergente
\end{frame}

%------------------------------------------------------------------------------%
\addtocounter{exemplo}{1}
\begin{frame}{Exemplo \theexemplo}
  Verifique a convergência da série
  \quad
  \(
    \sum_{n=1}^{\infty} a_n = \sum_{n=1}^{\infty} \frac{2n+1}{\;n^2+2n+1\;}
  \)

  \pause
  \vspace{\baselineskip}
  Comparamos no limite com a Série Harmônica
  \quad
  \(
    \sum_{n=1}^{\infty} b_n = \sum_{n=1}^{\infty} \frac{1}{\;n\;}
  \)

  \pause
  \vspace{0.3\baselineskip}
  \[
  \lim_{n\to\infty} \frac{a_n}{b_n}              \pause \;=\;
  \lim_{n\to\infty} \frac{(2n+1)n}{\;n^2+2n+1\;} \pause \;=\;
  \lim_{n\to\infty} \frac{4n+1}{\;2n+2\;}        \pause \;=\;
  \lim_{n\to\infty} \frac{\,4\,}{2}                     \;=\; 2
  \]

  \pause
  \vspace{0.5\baselineskip}
  Como a soma de \(b_n\) diverge então a de \(a_n\) também diverge
\end{frame}

%------------------------------------------------------------------------------%
\addtocounter{exemplo}{1}
\begin{frame}{Exemplo \theexemplo}
  Utilize o teste da comparação no limite para determinar se a série
  \[
    \sum_{n=2}^{\infty} \frac{1}{\,\ln n\,}
  \]
   converge ou diverge

   \pause
   \vspace{\baselineskip}
   Observe que $n=1$ não pode fazer parte da série pois $\ln 1 = 0$
\end{frame}

%------------------------------------------------------------------------------%
\begin{frame}{Exemplo \theexemplo}
  A série harmônica
  \[
    \sum_{n=2}^{\infty} \frac{1}{n}
  \]
  é divergente

  \pause
  \vspace{\baselineskip}
  Os termos da série original e da série harmônica são sempre positivos para $n\geq 2$.
\end{frame}

%------------------------------------------------------------------------------%
\begin{frame}{Exemplo \theexemplo}
  \onslide<+->{Aplicando o teste da comparação no limite}
  \begin{align*}
    \onslide<1->{\lim_{n\to\infty}\frac{a_n}{b_n}}
    \onslide<+->{& = \lim_{n\to\infty} \dfrac{\;\frac{1}{\ln n}\;}{\frac{1}{n}}         \\[2mm]}
    \onslide<+->{& = \lim_{n\to\infty} \frac{n}{\ln n}                                  \\[2mm]}
    \onslide<+->{& = \lim_{x\to\infty} \frac{x}{\ln x}                                  \\[2mm]}
    \onslide<+->{& = \lim_{x\to\infty} \frac{1}{\sfrac{1}{x}} & \text{Usando l'Hopital} \\[2mm]}
    \onslide<+->{& = \lim_{x\to\infty} x = \infty}
  \end{align*}
  \onslide<+->{Concluímos que a série diverge}
\end{frame}

% Exercício 24 página 31 Thomas?
%------------------------------------------------------------------------------%
\addtocounter{exemplo}{1}
\begin{frame}{Exemplo \theexemplo}
  A série
  \[
    \sum_{n=3}^{\infty} \frac{5n^3-3n}{n^2(n-2)(n^2+5)}
  \]
  converge ou diverge?
\end{frame}

%------------------------------------------------------------------------------%
\begin{frame}{Exemplo \theexemplo}
  Para $n$ suficientemente grande $a_n$ se comporta como
  \[
    b_n = \frac{n^3}{n^5} = \frac{1}{n^2}
  \]

  \pause
  \vspace{\baselineskip}
  A série $b_n$ é uma série $p$ com $p=2>1$ e portanto convergente

  \pause
  \vspace{\baselineskip}
  Além disso, $a_n$ e $b_n$ são maiores do que zero para todo $n$ suficientemente grande
\end{frame}

%------------------------------------------------------------------------------%
\begin{frame}{Exemplo \theexemplo}
  \begin{columns}
    \column{0.45\linewidth}
    \onslide<+->{Pelo teste da comparação no limite}
    \begin{align*}
        \onslide<1->{C & = \lim_{n\to\infty} \frac{a_n}{b_n}                                                 \\[2mm]}
        \onslide<+->{& = \lim_{n\to\infty} \dfrac{\; \dfrac{5n^3-3n}{\,n^2(n-2)(n^2+5)\,}\;}{\dfrac{1}{n^2}} \\[2mm]}
        \onslide<+->{& = \lim_{n\to\infty} \frac{5n^3-3n}{\;(n-2)(n^2+5)\;}                                  \\[2mm]}
        \onslide<+->{& = \lim_{n\to\infty} \frac{5n^3-3n}{\;n^3 -2n^2 +5n -10\;}                             }
      \end{align*}
    \column{0.55\linewidth}
    \begin{align*}
        \onslide<+->{C
        & = \lim_{n\to\infty}
            \frac{\dfrac{5n^3}{n^3}
                 -\dfrac{3n}{n^3}}
                {\;\dfrac{n^3}{n^3}
                 -\dfrac{2n^2}{n^3}
                 +\dfrac{5n}{n^3}
                  -\dfrac{10}{n^3}\;} \\[2mm]}
        \onslide<+->{& = \lim_{n\to\infty}
            \frac{5
                 -\dfrac{3}{n^2}}
                {\;1
                 -\dfrac{2}{n}
                 +\dfrac{5}{n^2}
                 -\dfrac{10}{n^3}\;} \\[2mm]}
        \onslide<+->{& = 5}
      \end{align*}
    \onslide<+->{Portanto a série original converge}
  \end{columns}
\end{frame}

%%------------------------------------------------------------------------------%
%\begin{frame}{Exemplo \theexemplo}
%  \begin{columns}
%    \column{0.5\linewidth}
%    Pelo teste da comparação no limite
%    \begin{align*}
%      C
%      & = \lim_{n\to\infty} \frac{a_n}{b_n}                                                   \\[2mm]
%      & = \lim_{n\to\infty} \dfrac{\; \dfrac{5n^3-3n}{\,n^2(n-2)(n^2+5)\,}\;}{\dfrac{1}{n^2}} \\[2mm]
%      & = \lim_{n\to\infty} \frac{5n^3-3n}{\;(n-2)(n^2+5)\;}                                  \\[2mm]
%      & = \lim_{n\to\infty} \frac{5n^3-3n}{\;n^3 -2n^2 +5n -10\;}                             \\[2mm]
%    \end{align*}
%    \column{0.5\linewidth}
%    Mudando o limite para uma variável real
%    \begin{align*}
%      C
%      & = \lim_{x\to\infty} \frac{5x^3-3x}{\;x^3 -2x^2 +5x -10\;}                             \\[2mm]
%      & = \lim_{x\to\infty} \frac{15x^2-3}{\;3x^2 -4x +5\;}                                   \\[2mm]
%      & = \lim_{x\to\infty} \frac{30x}{\;6x-4\;}                                              \\[2mm]
%      & = \lim_{x\to\infty} \frac{30}{6}                                                      \\[2mm]
%      & = 5
%    \end{align*}
%    Portanto a série original converge
%  \end{columns}
%\end{frame}

%------------------------------------------------------------------------------%
%------------------------------------------------------------------------------%
\section{Lista Mínima}

%------------------------------------------------------------------------------%
\begin{frame}{Lista Mínima}
  Estudar as Seção 6.5 da Apostila

  \vspace{\baselineskip}
  Exercícios: 4, 7

  \vfill
  Atenção: A prova é baseada no livro, não nas apresentações
\end{frame}

%------------------------------------------------------------------------------%
\end{document}
%------------------------------------------------------------------------------%
